\documentclass[12pt]{article}
\usepackage[utf8]{inputenc}
\usepackage[T1]{fontenc}
\usepackage[margin=1in]{geometry}
\usepackage{setspace}
\usepackage{xcolor}
\usepackage[colorlinks=true, urlcolor=blue!60!black, linkcolor=blue!60!black]{hyperref}
\setlength{\parindent}{2em}
\setlength{\parskip}{1.25em}
\renewcommand{\baselinestretch}{1.0}
\begin{document}
\begin{center}
{\LARGE \textit{Abrams v.\ United States}, 250 U.S.\ 616 (1919)\par}
\vspace{0.8em}
{\large Case Brief --- Four-Page Edition\par}
\end{center}
\vspace{1em}

\section{Memory Jogger}

Affirming the Espionage Act convictions of five Russian-born radicals who scattered leaflets urging a munitions strike, the Court held 7--2 that intent to hinder the war could be inferred from the likely effects of their speech; Justice Holmes's dissent, demanding a clear and imminent danger, became a foundation of modern free-speech law (\href{https://supreme.justia.com/cases/federal/us/250/616/}{Justia}).

\section{Detailed Case Facts}

In August 1918, during World War I, the United States sent expeditionary forces to Russia after the Revolution. Five Russian-born immigrants in New York City, Jacob Abrams, Hyman Lachowsky, Samuel Lipman, Mollie Steimer, and Hyman Rosansky, self-identified anarchists and a socialist, opposed the intervention. On a basement press they printed about 5,000 copies of two leaflets, distributing them secretly and throwing them from a window. The English leaflet, ``The Hypocrisy of the United States and her Allies,'' called President Wilson a ``coward'' and urged ``Workers of the World! Awake! Rise! Put down your enemy and mine!'' while repudiating German militarism. The Yiddish leaflet warned munitions workers that they were making ``bullets, bayonets, cannon, to murder not only the Germans, but also your dearest, best, who are in Russia,'' and declared that the reply ``has to be a general strike!''

\section{Procedural History}

The defendants were indicted on four conspiracy counts under the Espionage Act as amended in 1918: publishing ``disloyal, scurrilous and abusive language'' about the form of government; language intended to bring it into contempt; language intended to ``incite, provoke and encourage resistance to the United States in said war''; and urging curtailment of war production. A jury in the Southern District of New York convicted all five, with sentences of up to twenty years. They sought review by writ of error.

\section{Judicial Votes}

Affirmed, \textbf{7--2} (decided November 10, 1919).
\begin{itemize}
\item \textbf{Majority:} Justice Clarke, joined by Chief Justice White and Justices McKenna, Day, Van Devanter, Pitney, and McReynolds.
\item \textbf{Concurrences:} None.
\item \textbf{Dissent:} Justice Holmes, joined by Justice Brandeis.
\end{itemize}

\section{Holding}

Under \textit{Schenck} and \textit{Frohwerk}, the First Amendment did not shield the leaflets, and substantial evidence supported the convictions on Counts Three and Four. Applying a deferential bad-tendency and foreseeable-consequences approach, the Court held that speakers who adopt a plan that necessarily entails defeating the war effort against Germany are held to intend that result, whatever their pro-Russian motive.

\section{Analysis of Opinions}

\textbf{Majority (Clarke).} The First Amendment claim was ``definitely negatived'' by \textit{Schenck} and \textit{Frohwerk}. Because the sentences did not exceed what any one count allowed, the judgment stood if any count was supported. ``[M]en must be held to have intended, and to be accountable for, the effects which their acts were likely to produce.'' A general strike in munitions plants would halt production ``at the supreme crisis of the war'' in ``the greatest port of our land.'' Uncirculated writings, including ``Revolutionists! Unite for Action!,'' supplied ``much persuasive evidence.''

\textbf{Dissent (Holmes).} Eight months after writing \textit{Schenck}, \textit{Frohwerk}, and \textit{Debs}, Holmes recast ``clear and present danger'' as a speech-protective rule. Statutorily, attacking the President's policy is not attacking the ``form of government,'' Count Three required ``some forcible act of opposition,'' and Count Four required that hindering the war be the actor's ``proximate motive'': otherwise a patriot urging less airplane spending could be jailed. Constitutionally, Congress may punish speech only when it produces or is intended to produce a ``clear and imminent danger.'' The ``surreptitious publishing of a silly leaflet by an unknown man'' posed no such threat, and the sentences punished the defendants ``for the creed that they avow.'' Rejecting seditious libel, Holmes argued that because ``time has upset many fighting faiths,'' the Constitution bets on ``free trade in ideas.''

\textbf{Later.} \textit{Brandenburg v.\ Ohio} (1969) vindicated Holmes: advocacy may be punished only if ``directed to inciting or producing imminent lawless action and is likely to incite or produce such action.''

\section{Examples: Future Applications}

\begin{enumerate}
\item \textbf{Same-Side: Material Support.} Speech coordinated with a designated foreign terrorist organization may be punished, as in \textit{Holder v.\ Humanitarian Law Project} (2010), echoing the majority's wartime deference.
\item \textbf{Same-Side: Storming a Recruiting Station.} A speaker telling a crowd to break in now and burn draft files, which it does, can be convicted even under Holmes's test and \textit{Brandenburg}.
\item \textbf{Opposite-Side: A Strike Next Month.} Urging dockworkers to strike next month against arms shipments is not directed to imminent lawless action.
\item \textbf{Opposite-Side: Calling the President a Coward.} Prosecuting ``disloyal, scurrilous and abusive language'' revives the seditious libel that \textit{New York Times v.\ Sullivan} condemned.
\item \textbf{Fence-Sitter: A Viral Call to Shut Down a Plant.} An anonymous post urging followers to ``shut down'' a defense plant ``right now,'' followed by an armed crowd within twenty minutes, may meet imminence, but intent is debatable.
\end{enumerate}

\section{Critique}

\textbf{Holmes's unacknowledged pivot.} Holmes claimed he ``never [had] seen any reason to doubt'' \textit{Schenck}, \textit{Frohwerk}, and \textit{Debs}, yet his dissent quietly turned ``clear and present danger'' from a deferential wartime attempt test into a strict imminence and specific-intent rule.

\textbf{Limits of the market metaphor.} Critics note that ``free trade in ideas'' assumes rational listeners, cheap verification, and symmetric access, assumptions strained by cognitive bias, concentrated media power, and viral disinformation.

\section{Key Quotations}

\begin{quote}
``But when men have realized that time has upset many fighting faiths, they may come to believe even more than they believe the very foundations of their own conduct that the ultimate good desired is better reached by free trade in ideas -- that the best test of truth is the power of the thought to get itself accepted in the competition of the market, and that truth is the only ground upon which their wishes safely can be carried out. That, at any rate, is the theory of our Constitution. It is an experiment, as all life is an experiment.''
\end{quote}
\noindent --- Justice Holmes, dissenting, 250 U.S.\ at 630.

\begin{quote}
``While that experiment is part of our system, I think that we should be eternally vigilant against attempts to check the expression of opinions that we loathe and believe to be fraught with death, unless they so imminently threaten immediate interference with the lawful and pressing purposes of the law that an immediate check is required to save the country.''
\end{quote}
\noindent --- Justice Holmes, dissenting, 250 U.S.\ at 630.

\begin{quote}
``It is only the present danger of immediate evil or an intent to bring it about that warrants Congress in setting a limit to the expression of opinion where private rights are not concerned.''
\end{quote}
\noindent --- Justice Holmes, dissenting, 250 U.S.\ at 628.

\end{document}
